% Abstract (Task 42): central claim of handoff/37_narrative_v2.md section 1.1 with the A6x result.
\begin{abstract}
Attention weights of transformer-based network intrusion detection systems (NIDS) are often read as explanations, and
agreement between attention and attribution methods has been proposed as an attack signal.
We tested both readings on FT-Transformers (feature tokenizer transformers; \nSeeds{} seeds) trained on LYCOS-IDS2017
and LycoS-Unicas-IDS2018, using their \nCommonClasses{} common attack classes and excluding evaluation rows with a near
duplicate in the training set.
Rollout attention agreed weakly to moderately with attribution methods (Spearman correlation with GradientSHAP
\spRolloutGShapVoneSeventeen{} on 2017 and \spRolloutGShapVoneEighteen{} on 2018), was less faithful than Integrated
Gradients and occlusion, and was stiff under protocol-constrained attacks that targeted it.
Attention-attribution consistency scores added to a confidence-and-distance detector helped against an attacker that
ignored them, as did controls without attention.
Against attackers that optimised the exact score of the detector they faced, with surrogate gradients only for the
\nRankScores{} rank-based scores, the consistency scores instead raised the share of successful and undetected attack
rows, by \axHeadVoneSeventeenAbs{} on 2017 (\paramCILevel{} confidence interval \axHeadVoneSeventeenCI) and by
\axHeadVoneEighteenAbs{} on 2018 (\axHeadVoneEighteenCI), with \axHeadOppSeedsPhrase{} of the
opposite sign.
An attacker with differentiable stand-ins for the scores had shown a reduction of \axStandinVoneSeventeenAbs{} on 2017;
giving the attacker of the base detector the same iterate selection as the exact-objective attacker changed no row.
Adaptive evaluations of explanation-based attack detectors need attackers that optimise the detector's exact score.
\end{abstract}
